\title{Byte Latent Transformer: Patches Scale Better Than Tokens}

\begin{abstract}
We present the Byte Latent Transformer ({\textsc BLT}), a novel large language model (LLM) framework that operates directly at the byte level. {\textsc BLT} achieves parity with traditional tokenization-dependent LLMs in terms of performance, while delivering notable gains in inference speed and resilience. The {\textsc BLT} architecture represents sequences as patches composed of bytes with variable lengths, which are used as its fundamental computational elements. Patch boundaries are adaptively determined according to the predicted entropy for the next byte, assigning greater computation and model resources to regions with higher informational content. We conduct, for the first time, a scaling analysis of byte-level models controlled for floating point operations (flops), training systems with up to 8B parameters on 4T bytes. These findings establish the practicality of scaling architectures trained on unprocessed byte streams without the need for a predefined vocabulary. Training and inference processes become more efficient because the model opts for longer patches in more predictable regions of the data. Additionally, qualitative advances are observed in complex reasoning tasks and in handling rare data. At any given inference budget, {\textsc BLT} demonstrates superior scaling capabilities compared to tokenization-dependent models, enabled by concurrent expansion of both the patch and model dimensions.
\end{abstract}

\section{Introduction}
\label{section:intro}

We present the Byte Latent Transformer~(\textbf{BLT}), an architecture that forgoes the need for tokenization by directly processing raw byte sequences. Remarkably, BLT is the first architecture of its kind to attain parity in performance with conventional tokenization-based models at scale, while delivering marked gains in both computational efficiency and robustness (\S\ref{sec:robustness}). In standard practice, large language models (llms) are trained in an almost completely end-to-end fashion, with the notable exception of tokenization---a heuristic pre-processing mechanism that aggregates bytes into a predetermined set of discrete tokens. This process introduces representational biases in the encoding of text, which result in various limitations: sensitivity to domain or modality shifts~\citep{dagangetting}, heightened vulnerability to noisy inputs~(\S\ref{sec:robustness}), insufficient modeling of orthographic structure~\citep{edman2024cute}, and inequitable treatment across languages~\citep{liang2023xlm,petrov2024language,limisiewicz2024myte}.

In the past, tokenization has been a necessary step because training llms directly on raw bytes is computationally prohibitive at large scales due to extended sequence lengths~\citep{xue2022byt5}. Earlier research addresses this challenge by adopting more resource-efficient self-attention methods~\citep{el2020characterbert, clark2022canine} or exploring architectures that eliminate attention altogether~\citep{wang2024mambabyte}~(\S\ref{section:related_work}). Nonetheless, these strategies are particularly effective for \textit{small models}. When scaling up, the bulk of the computation in a Transformer comes from the sizable feed-forward network layers applied to each byte, rather than from the attention module itself.

To optimize the use of compute resources, we introduce an adaptive, trainable approach for aggregating bytes into \textit{patches}~(\S\ref{section:patching}), alongside a novel architecture that integrates information from both individual bytes and their aggregated patches. In contrast to tokenization, {\textsc BLT} does not rely on a pre-defined set of patches. Instead, it employs compact, learned encoder and decoder modules that transform arbitrary byte groupings into latent patch representations. Our findings demonstrate that this strategy achieves \textit{greater} efficiency in compute allocation compared to models based on tokenization.

Tokenization-based LLMs uniformly distribute computational resources across all tokens. This approach prioritizes consistency over optimal performance, as token boundaries—shaped by compression strategies—do not always reflect the actual difficulty of prediction tasks. Our model, by contrast, is founded on the principle that computational effort should be directed adaptively to those regions where it is most required. For instance, predicting word endings is generally a straightforward, low-entropy task and does not necessitate the full capacity of a large transformer, whereas selecting the initial word in a sentence demands a more substantial allocation of resources. In line with this principle, the {\textsc BLT} architecture~(\S\ref{section:architecture}) employs a modular design featuring three transformer components: a pair of compact byte-level \textit{local models} and a single, more powerful global \textit{latent transformer}~(\autoref{fig:arch_high_level}).
To facilitate efficient computation, {\textsc BLT} organizes input data into patches based on the entropy associated with next-byte prediction, thereby producing context-dependent groups of bytes that maintain similar levels of information density.

We introduce the first comprehensive study on scaling byte-level models under fixed computational budgets, exploring models up to 8B parameters and 4T training bytes. Our results demonstrate that it is feasible to train models end-to-end directly from bytes, bypassing the need for predetermined token vocabularies. In terms of training efficiency, {\textsc BLT} achieves parity in flop-controlled performance\footnote{Model computational requirements are quantified based on the count of Floating Point OPerations (flops) necessary.} with Llama 3, while reducing inference-time flop consumption by as much as 50\% (\S\ref{section:scaling}).

Additionally, we find that operating directly on raw bytes substantially enhances modeling of rare or infrequent data features. Compared to models relying on tokenizers, {\textsc BLT} demonstrates increased resilience to noisy inputs and exhibits superior capability in tasks such as orthographic processing, phonological generalization, and translation for low-resource languages (\S\ref{sec:robustness}).

Moreover, with the {\textsc BLT} architecture, it becomes possible to simultaneously scale both the model's parameter count and patch length without exceeding a fixed inference flop limit. By increasing patch sizes, computational savings accrue, allowing further growth of the global latent transformer, which operates less frequently as a consequence. Our inference-flop-constrained scaling experiments (\autoref{fig:fixed_inference_scaling}) reveal that byte-level models show substantially improved scaling behavior relative to models based on fixed-tokenization approaches.

In conclusion, this work offers several key contributions: 1) We present {\textsc BLT}, a novel byte latent LLM architecture capable of adaptively allocating computational resources to enhance flop efficiency, 2) We provide evidence that {\textsc BLT} can reach parity with Llama 3 in terms of flop-controlled training up to an 8B parameter scale, while also permitting significant increases in flop efficiency—up to 50\%—by sacrificing minimal accuracy in evaluation metrics, 3) Our {\textsc BLT} models enable an alternative paradigm for scaling LLMs, permitting expansion of model size without increasing inference costs, 4) We empirically verify that {\textsc BLT} demonstrates superior resilience to input noise and possesses an improved sensitivity to sub-word patterns in the input compared to conventional token-based LLMs. The training and inference code for {\textsc BLT} is publicly available at~\url{https://github.com/facebookresearch/blt}.

\section{Patching: From Individual Bytes to Groups of Bytes}
\label{section:patching}

Dividing byte sequences into discrete \textit{patches} enables {\textsc BLT} to flexibly distribute computational resources according to the surrounding context. As illustrated in Figure~\ref{fig:patching_types}, there are various approaches for splitting bytes into patches. More specifically, a patching function $f_p$ partitions a byte sequence $\pmb{x}=\{x_i,|i=1,\ldots n\}$ of length $n$ into $m < n$ patches, denoted $\pmb{p}=\{p_j|j=1,\ldots,m\}$, by associating each $x_i$ with a value in \{0,1\}, where a value of 1 signifies the beginning of a new patch. In both token-centric and patch-based architectures, the main determinant of computational requirements is the number of main Transformer steps. For {\textsc BLT}, this corresponds to the total number of patches used to represent the input given a particular patching strategy. Thus, the key determinant of both training and inference efficiency with a given patching method is the mean number of bytes per patch, referred to as the \textit{patch size}~(\S\ref{section:flops}).

We proceed by detailing three forms of patching functions: segmentation into patches with a constant byte count~(\S\ref{section:static-patch}), patching at whitespace boundaries~(\S\ref{section:white-patch}), and adaptively defining patches based on the entropy estimates from a lightweight byte language model~(\S\ref{section:dyn-patch}). Subsequently, we explore the concept of incremental patching and clarify distinctions between patching and conventional tokenization methods~(\S\ref{section:bpe-patch}).

\subsection{Strided Patching Every K Bytes}
\label{section:static-patch}

One common approach to organizing bytes is to divide them into fixed-length patches of size $k$, as seen in MegaByte~\citep{yu2023megabyte}. This method, which relies on a constant stride, is simple to both implement and use during model training and inference. It offers an intuitive means to adjust the average patch size, thereby facilitating precise control over computational overhead (FLOPs). Despite these practical benefits, fixed-size patching is accompanied by notable drawbacks. Primarily, computational resources are not allocated according to the content's complexity or relevance: for instance, a transformer's step $j$ may be underutilized if it merely predicts whitespace within a code snippet, or inadequately applied when processing highly information-rich segments, such as those containing mathematics. Additionally, this technique causes variability and a lack of contextual awareness in patching, so that even identical byte sequences—such as recurring words—may be segmented differently each time.

\subsection{Space Patching}
\label{section:white-patch}

\citet{slagle2024spacebyte} introduces a straightforward yet impactful modification to the standard strided patching approach by generating new patches following any space-like byte\footnote{
    Space-like bytes refer to any byte that is neither a Latin character, a numeral, nor a utf-8 continuation byte. Furthermore, patches are required to contain at least one byte that is not space-like. }—positions which often correspond to natural linguistic boundaries across a variety of languages. In the Space patching approach, each word is given a dedicated latent transformer step (i.e., additional computational effort), allowing the model to allocate more flops to the modeling of every individual word. This results in consistent patching of words within sequences and ensures computational resources are dedicated to predictions that are challenging, which frequently occur at space boundaries. For instance, generating the initial byte of a response to ``Who composed the Magic Flute?\uline{\hspace{.6em}}'' demands considerably more effort compared to the subsequent bytes after ``M,'' as the initial character narrows the pool of possible completions—making it much easier to produce ``Mozart'' after the initial ``M'' is predicted. Nonetheless, the Space patching method is limited in that it cannot seamlessly accommodate all languages and domains, and—crucially—lacks the ability to flexibly adjust patch sizes. In what follows, we describe a novel patching strategy inspired by the observation that initial bytes of words are generally the most challenging to predict, offering an intuitive means of varying patch size.

\subsection{Entropy Patching: Using Next-Byte Entropies from a Small Byte LM}
\label{section:dyn-patch}

Instead of depending on traditional heuristics like whitespace, our method adopts a data-driven strategy to detect next-byte predictions with elevated uncertainty. We propose \textit{entropy patching}, a technique that employs entropy measurements to determine the locations of patch boundaries.

We train a small byte-level auto-regressive language model on the training data for {\textsc BLT} and compute next byte entropies under the LM distribution $p_{e}$ over the byte vocabulary $\mathcal{V}$:

\begin{equation}
H(x_i) = \sum_{v \in \mathcal{V}} p_{e}(x_i=v|\pmb{x}_{<i}) \log p_{e}(x_i=v|\pmb{x}_{<i})
\end{equation}

We investigate two approaches for detecting patch boundaries using entropies $H(x_i)$. The initial method selects points exceeding a global entropy threshold, as shown in~\autoref{fig:patching}. The alternative technique selects points that display a significant increase compared to the preceding entropy value. This latter method can additionally be viewed as pinpointing locations where the approximately monotonically decreasing entropy within a patch is disrupted.

\begin{align*}
    \text{Global Constraint}\hspace{1cm}H(x_t)& > \theta_g\\
    \text{Approx. Monotonic Constraint}\hspace{1cm}H(x_t)& - H(x_{t-1}) > \theta_r
\end{align*}

Patch boundaries are determined through a simple preprocessing procedure performed while loading the data. This approach contrasts with that of~\citet{nawrot-etal-2023-efficient}, who train a classifier to detect entropy-based patch boundaries. In our experimental evaluation~(\S\ref{section:experiments}), we examine and compare these two techniques for discriminating between bytes with low and high entropy.

\subsection{The Byte-Pair Encoding (BPE) Tokenizer and Incremental Patching}
\label{section:incremental-patching}
\label{section:bpe-patch}

A large number of contemporary llms—including the baseline Llama 3—implement subword tokenization schemes such as bpe~\citep{gage1994new, sennrich-etal-2016-neural}. Throughout this work, the term ``tokens'' designates byte-groupings selected from a \textit{finite} vocabulary that is fixed before the training process, in contrast to ``patches,'' which are defined as flexible, dynamically assembled byte sequences lacking a predetermined vocabulary. One notable distinction between patches and tokens lies in the model’s inability to access the original byte-level features when using tokens.

An important advancement introduced by {\textsc BLT} compared to models that rely on tokenization is its reconceptualization of the balance between vocabulary size and computational requirements. In conventional large language models, expanding the vocabulary typically results in tokens that encompass more characters on average, which reduces the total number of prediction steps needed; however, this also necessitates a larger output dimension in the model's final projection layer. Consequently, this compromise restricts the extent to which tokenization-based methods can vary token sizes or minimize inference overhead. As an illustration, Llama 3 achieves an increase in average token length—from 3.7 to 4.4 bytes—by enlarging its embedding matrix by a factor of four relative to Llama 2.

During generation, {\textsc BLT} must determine at each position in the byte sequence whether it is at a patch boundary, since this influences if additional computation through the Latent Transformer is required. This choice must be made solely based on information available up to the current position; it cannot rely on bytes that have not yet been produced. Therefore, any patching approach must partition the byte sequence without knowledge of subsequent bytes. More precisely, a patching method $f_p$ possesses the property of incremental patching if the following holds:
$$f_p(\pmb{x}_{<i}) = f_p(\pmb{x})_{<i}$$
bpe does not constitute an incremental patching scheme, as a given prefix could result in different tokenizations depending on its suffix, and thus it fails to meet the criterion stated above\footnote{One can modify bpe to be incremental by adding a special delimiter to mark patch boundaries, but this leads to a longer byte sequence.}.

\section{{\textsc BLT} Architecture}
\label{section:architecture}
{\textsc BLT} integrates a sizeable global autoregressive language model tasked with processing patch-level representations, supplemented by two lightweight local modules that respectively transform byte sequences into patches and reconstruct bytes from patch representations~(\autoref{fig:arch_high_level}).

\subsection{Latent Global Transformer Model}

\textit{The Latent Global Transformer} refers to an autoregressive transformer architecture, denoted as $\mathcal{G}$, comprised of $l_{\mathcal{G}}$ layers. This model accepts as input a sequence of latent patch representations $p_j$ and produces a corresponding sequence of output patch representations $o_j$. In this work, patch indices are represented by the subscript $j$, while byte indices use $i$. The global transformer employs a block-causal attention mask~\citep{dubey2024llama}, ensuring that each token can attend only to tokens within and prior to the current patch in the same document. Notably, the global model is responsible for the majority of computational cost during both pre-training and inference. Therefore, by modulating when the global transformer is utilized, one can directly influence and modulate the total computation allocated to different segments of the input and output depending on their complexity.

\subsection{Local Encoder}
\label{section:local}

\textit{The Local Encoder Model} $\mathcal{E}$ is designed as a compact transformer variant, comprising only $l_{\mathcal{E}} << l_{\mathcal{G}}$ layers. Its principal objective is to efficiently translate sequences of input bytes $b_i$ into high-quality patch representations $p_j$. Distinguishing it from standard transformer architectures, this model introduces a cross-attention mechanism placed after each transformer layer, whose purpose is to aggregate the byte-level information into patch-level representations~(\autoref{fig:crossattn}).

Initially, the sequence of input bytes $b_i$ is projected via an embedding matrix in $\mathbb{R}^{256 \times h_{\mathcal{E}}}$, yielding $x_i$. These token embeddings may be further supplemented by hash-embeddings as an optional enhancement~(\S\ref{sec:hash-emb}).

The subsequent stages involve iteratively applying transformer layers followed by cross-attention modules~(\S\ref{sec:enc-cross-attn}), progressively molding the representations into patch vectors $p_i$ to be utilized by the overarching global transformer $\mathcal{G}$.

Within each transformer layer, a \textit{local block causal} attention mask is imposed: each byte can attend to the $w_{\mathcal{E}}$ tokens immediately before it, potentially spanning multiple patch boundaries as they dynamically form, yet strictly observing document limits. The subsections below elaborate on the embedding strategies and the specifics of the cross-attention mechanism.

\subsubsection{Encoder Hash n-gram Embeddings}
\label{sec:ngram-emb}
\label{sec:hash-emb}
An essential factor in generating resilient and informative representations at every position $i$ is including context from previous bytes. Within {\textsc BLT}, we address this by capturing information about the current byte $b_i$ in isolation as well as within the scope of a byte n-gram. Specifically, at each position $i$, we begin by forming byte-grams

\begin{align}
    g_{i,n}=\{b_{i-n + 1},\ldots, b_i\}
\end{align}

for every byte position $i$ and for values of $n$ ranging from three to eight.\footnote{
Byte-grams with length $n$ or greater are excluded if $i < n$. }

Following this, we define hash $n$-gram embeddings, which assign all byte $n$-grams to indices in an embedding table $E_{n}^{hash}$ of a predetermined size for each $n$ in $\{3,4,5,6,7,8\}$ by means of a hash function~\citep{bai2010learning}. The generated embedding is summed with the embedding of the corresponding byte, after which the combined vector is normalized and supplied to the local encoder. The computation for the enriched embedding is given by

\begin{align}
e_i &= x_i + \sum_{n = 3,...,8}E_{n}^{hash}(\text{Hash}(g_{i,n})) \\
\text{where, Hash}(g_{i,n}) &= \text{RollPolyHash}(g_{i,n}) \% |E_{n}^{hash}|
\end{align}

We scale $e_i$ by dividing it by the total count of $n$-gram sizes increased by one, and employ the RollPolyHash technique as described in Appendix~\ref{appendix:rollpolyhash}. Section~\ref{section:ablations} presents an analysis that isolates the impact of $n$-gram hash embeddings, considering various $n$ values and different embedding table dimensions, on scaling law dynamics under fixed computation. Alongside hash-based $n$-gram embeddings, we conducted trials using frequency-based $n$-gram embeddings; further specifics on these experiments are available in Appendix~\ref{appendix:ngrams}.

\subsubsection{Encoder Multi-Headed Cross-Attention}
\label{sec:enc-cross-attn}
Our approach is largely inspired by the input cross-attention mechanism used in the Perceiver architecture~\citep{jaegle2021perceiver}, with the primary distinction that our latent representations are derived from the variable patch representations rather than from a predefined set of fixed latent vectors~(\autoref{fig:crossattn}); moreover, these representations attend exclusively to the bytes constituting their associated patch. Within this module, each patch $p_j$ is paired with a corresponding query vector, created by aggregating the byte representations of $p_j$ and subsequently applying a linear transformation, $\mathcal{E}_{C} \in \mathbb{R}^{h_{\mathcal{E}} \times (h_{\mathcal{E}} \times U_{\mathcal{E}})}$, where $U_{\mathcal{E}}$ specifies the number of cross-attention heads in the encoder. More precisely, denoting the byte sequence of patch $p_j$ as $f_{\text{bytes}}(p_j)$, we compute

\begin{align}
P_{0,j} &= \mathcal{E}_{C}(f_\text{bytes}((p_j)), f~ \text{is a pooling function} \\
P_l &= P_{l-1} + W_o\left(\text{softmax}\left(\frac{QK^T}{\sqrt{d_k}}\right)V\right) \\
\text{where } Q_j &= W_q(P_{l-1,j}), K_i = W_k(h_{l-1,i}), V_i = W_v(h_{l-1,i}) \\
h_l &= \text{Encoder-Transformer-Layer}_l(h_{l-1})
\end{align}

Here, $P \in \mathbb{R}^{n_p \times h_{\mathcal{G}}}$ denotes $n_p$ patch representations intended for the global model. These representations are constructed by aggregating the byte-level embeddings $e_i$ associated with each individual patch $p_j$. The matrices $W_q$, $W_k$, $W_v$, and $W_o$ denote the linear projections for the queries, keys, values, and outputs, with the keys and values specifically derived by projecting the byte-level features $h_i$ from the preceding layer (using $e_i$ as input for the initial layer). A targeted masking scheme is applied during patching such that a given query $Q_j$ only has access to keys and values originating from the bytes that comprise patch $j$. As multi-headed attention is utilized on $Q$, $K$, and $V$, and given that patch representations generally have higher dimensionality ($h_{\mathcal{G}}$) relative to $h_{\mathcal{E}}$, we represent $P_l$ across multiple attention heads of size $h_{\mathcal{E}}$ for the cross-attention computation and subsequently concatenate these outputs to restore a total dimensionality of $h_{\mathcal{G}}$. Further, a pre-LayerNorm is applied to each of the queries, keys, and values before attention, with positional embeddings deliberately omitted within this cross-attention design. The output from the cross-attention block is combined with its input via a residual (skip) connection.

\subsection{Local Decoder} 

Analogous to the local encoder, the local decoder $\mathcal{D}$ consists of a compact transformer architecture with $l_{\mathcal{D}} << l_{\mathcal{G}}$ layers. Its role is to transform a sequence of global patch representations $o_j$ into the original byte sequence $y_i$. The decoder generates each byte conditioned on the sequence of bytes it has already produced, relying on the contextual hidden states obtained from the local encoder corresponding to the byte sequence. It processes information through $l_{\mathcal{D}}$ repeated pairs of cross-attention modules and transformer blocks. In this setup, the cross-attention module precedes the transformer block within each layer, allowing the model to first map patch-level representations into byte-level representations, after which the subsequent transformer block refines these representations within the byte sequence.

\subsubsection{Decoder Multi-headed Cross-Attention} 

In the decoder cross-attention, the roles of the queries and key/values are interchanged i.e. the byte-representations are now the queries, and the patch representations are now the key/values. The initial byte-representations for the cross-attention are initialized as the byte embeddings from the last encoder layer i.e. $h_{l_{\mathcal{E}}}$. The subsequent byte-representations for layer $l$, $d_{l,i}$ are computed as:

\begin{align}
D_0 &= h_{l_{\mathcal{E}}} \\
B_l &= D_{l-1} + W_o\left(\text{softmax}\left(\frac{QK^T}{\sqrt{d_k}}\right)V\right), \\
  &\text{where } Q_i = W_q(d_{l-1,i}), K_i = W_k(\mathcal{D}_C(o_j)), V_i = W_v(\mathcal{D}_C(o_j)) \\
  D_l &= \text{Decoder-Transformer-layer}_l(B_l)
\end{align}

Here, $W_k$ and $W_v$ denote the key and value projection matrices, which perform linear transformations and splitting, denoted by $\mathcal{D}_C$, applied to the final patch representations $o_j$ obtained from the global model. $W_q$ serves as the query projection matrix and is applied to the byte-level representations $d_{l-1}$ from the preceding decoder transformer layer (or to $h_{l_{\mathcal{E}}}$ for the initial layer). The output projection is carried out by $W_o$, resulting in $B \in \mathbb{R}^{h_{\mathcal{D}} \times n_b}$, where $n_b$ corresponds to the number of output bytes. The subsequent decoder representations $D_l$ are produced by applying a decoder transformer layer to the output $B$ of the cross-attention component. Following the approach of the local encoder’s cross-attention mechanism, we adopt multi-head attention, incorporate pre-LayerNorm, omit positional embeddings, and introduce a residual connection around the cross-attention block.

\section{Experimental Setup}
\label{section:experiments}
We establish a rigorous experimental protocol aimed at fairly comparing {\textsc BLT} against models based on tokenization, specifically ensuring that {\textsc BLT} does not benefit from potential advantages associated with accessing longer sequence contexts.

\subsection{Pre-training Datasets} All model variants explored in this study are initially trained on two main datasets: 1) The Llama 2 dataset~\citep{touvron2023llama}, containing 2 trillion tokens sourced from an assortment of public repositories and meticulously processed to enhance data quality; and 2) {\textsc BLT}-1T, a newly constructed collection consisting of 1 trillion tokens sourced from a variety of publicly accessible datasets, including a segment of pre-training material released by Datacomp-LM~\citep{li2024datacomp}. The former serves as the basis for scaling law analyses, particularly those involving the optimal token count as identified in~\cite{dubey2024llama}, and informs architectural optimizations for {\textsc BLT}. The latter is employed in a comprehensive pre-training phase in order to provide a fair comparison to Llama 3 in subsequent evaluation tasks. Both datasets deliberately exclude any content derived from Meta platforms or services. Additionally, for benchmarking with tokenizer-based models, we employ the Llama 3 tokenizer featuring a vocabulary of 128K tokens, as it delivered superior baseline outcomes compared to the Llama 2 tokenizer within our testing framework.

\subsection{Entropy Model} In our study, the entropy model is implemented as a byte-level language model, trained on the same data distribution as the full {\textsc BLT} model. By default, we utilize a transformer architecture consisting of 100M parameters, 14 layers, and a hidden size of 512, employing sliding window attention that spans 512 bytes. All other hyperparameter choices align with those of our local and global transformer models. Further, as described in \autoref{section:ablations}, we have explored varying the model size, receptive field, and architecture. Notably, for sufficiently small receptive fields, the learned entropy model can be represented effectively as a compact lookup table.

\subsection{Entropy Threshold and Equalizing Context Length} For models utilizing entropy-based patching, we determine a patching threshold to produce a targeted average \textit{patch size} on the pretraining data mixture. In the {\textsc BLT} setting, the \textit{patch size} is not constrained by tokenization and can be freely selected, which can considerably influence the amount of context available to the model. To ensure fairness and prevent models with larger patch sizes from benefiting from increased context, we standardize the expected number of bytes processed in each batch. Consequently, when increasing the patch size, we proportionally decrease the sequence length. Specifically, with Llama 2 data, we adopt an 8k byte context, whereas on the {\textsc BLT}-1T dataset, we expand the context to 16k bytes on average, keeping the average batch size fixed at 16M bytes.

Although the mean batch size remains fixed, dynamic patching approaches result in varying proportions of bytes per patch when assembling data batches. To maximize efficiency, our approach to {\textsc BLT} training involves grouping batches by patches, thereby circumventing the need for padding within the computationally intensive latent transformer. This design choice guarantees a uniform patch count across all batches. To further control memory usage, byte sequences are padded or, if necessary, truncated to lengths of 12k and 24k bytes for the Llama 2 and {\textsc BLT}-1T datasets, respectively, minimizing the impact of exceptionally large patches on memory consumption during training.

\subsection{Entropy Model Context}
\label{section:entropy-context}
Our empirical observations indicate that entropy patching tends to produce increasingly larger patches within highly structured inputs, such as multiple choice questions, which are typically marked by significant repetition (see the example of patching applied to an MMLU item in \autoref{fig:mmlu_patching}). These changes are attributable to decreased entropy in the duplicated material present in the entropy model context. Consequently, for the large-scale experiment with {\textsc BLT}-Entropy employing a patch size of 4.5, we reinitialize the entropy context with line breaks and utilize the approximate monotonicity constraint, as this approach better mitigates the issue of "entropy drift" stemming from variations in context length. While this adjustment impacts only the manner in which entropy values are calculated, our method for determining the entropy threshold remains unchanged.

\subsection{FLOPs Estimation} 
\label{section:flops}

Our calculation of transformer flops is primarily based on the methodology presented in Chinchilla~\citep{hoffmann2022training}, including the flops associated with the feed-forward components, the qkvo projections in self-attention, as well as the flops required for attention computation and output projection. One key distinction in our approach is the assumption that the input embedding layer utilizes an efficient lookup table, in contrast to a dense matrix multiplication; as a result, we treat this as incurring zero flops. Consistent with established work, we also consider the backwards pass to require twice the flop count of the corresponding forward pass.

To compute flops \textit{per byte} for {\textsc BLT} models, we add up the flops for the local encoder transformer, the global latent transformer, and the local decoder transformer, together with the cross attention blocks in the encoder and the decoder:

\begin{align}
    \text{FL}_{\text{{\textsc BLT}}} &= \frac{1}{n_p}\,\text{Transf. FL}(h_{\mathcal{G}}, l_{\mathcal{G}}, m=n_{ctx}/n_p, V=0) \\
   &+ \text{Transf. FL}(h_{\mathcal{E}}, l_{\mathcal{E}}, m=w_{\mathcal{E}}, V=0) \\
   &+ \text{Transf. FL}(h_{\mathcal{D}}, l_{\mathcal{D}}, m=w_{\mathcal{D}}, V=256) \\ 
    &+ \frac{k}{n_p} \, \text{Cross Attn. FL}(h_{\mathcal{E}}, l_{\mathcal{E}}, m=n_p, r=n_p/k) \\ 
   &+ \text{Cross Attn. FL}(h_{\mathcal{D}}, l_{\mathcal{D}}, m=k, r=k/n_p)
\end{align}

Here, $n_{ctx}$ refers to the input sequence length measured in bytes, while $n_p$ denotes the size of each patch. The parameter $r$ signifies the proportion of queries relative to keys and values, and $k$ defines the ratio between the dimensions of patches and bytes; this indicates how many local partitions of the model are combined to create a unified global representation ($k=2$ in \autoref{fig:crossattn}). $V$ represents the size of the vocabulary used exclusively by the local decoder’s output projection. The choice between applying a module on either the byte sequence or the patch sequence determines the attention context length, $m$. Consequently, we adapt the flop calculations for each specific component. The precise formulas used for calculating flops in Transformer-FLOPs and Cross-Attention FLOPs are included in Appendix~\ref{section:flops-appendix}.

\subsection{Bits-Per-Byte Estimation} Perplexity only makes sense in the context of a fixed tokenizer as it is a measure of the uncertainty for each token. When comparing byte and token-level models, following previous work~\citep{xue2022byt5, yu2023megabyte, wang2024mambabyte}, we instead report Bits-Per-Byte (BPB), a tokenizer independent version of perplexity. Specifically:

\begin{align}
    \text{BPB}(x)&= \frac{\mathcal{L}_{CE}(\pmb{x})}{\ln(2)\cdot n_{\text{bytes}}}
\end{align}

In this formulation, the cumulative cross-entropy loss, which quantifies the uncertainty associated with the data $\pmb{x}$, is adjusted by dividing by both the total byte count of $\pmb{x}$ and a normalization constant.

\subsection{Transformer Architecture Hyperparameters} In configuring all transformer blocks within {\textsc BLT}, encompassing both local and global components, our approach generally aligns with the architectural choices made in Llama~3~\citep{dubey2024llama}. The feed-forward modules incorporate the SwiGLU activation~\citep{swiglu}, while rotary positional embeddings (RoPE)~\citep{RoPe} with a setting of $\theta=500000$~\citep{xiong2024effective} are exclusively applied within the self-attention layers. For normalization across layers, we rely on RMSNorm~\citep{RMSNorm}. All self-attention blocks employing attention masks such as \textit{block causal} or \textit{fixed-window block causal} utilize Flash attention~\citep{dao2022flashattention}, and the window size for fixed-width attention is established at 512. As the cross-attention modules necessitate dynamically determined, patch-specific masks, we leverage Flex Attention\footnote{\url{https://pytorch.org/blog/flexattention}}, allowing us to take advantage of fused kernels that notably accelerate the training process.

\subsection{{\textsc BLT}-Specific Hyperparameters} To evaluate the performance of {\textsc BLT} models, our experimental design encompasses two main avenues: analyzing scaling behavior and performing downstream task assessments, with models spanning various sizes (400M, 1B, 2B, 4B, and 8B parameters). Architectural hyperparameter configurations for these different model scales are detailed in Appendix~Table~\ref{tab:arch}. For the local encoder's initial cross-attention layer, queries are initialized using max-pooling. Across all {\textsc BLT} variants, we utilize $500,000$ hashes computed via a single hash function, considering n-grams ranging in size from 3 to 8. The learning rate is set at $4e-4$ uniformly for all model sizes. This consistent learning rate was selected based on hyperparameter tuning over the range $1e-3$ to $1e-4$ at the 400M and 1B scales, revealing no advantage to further adjustment. Training batch sizes for Llama-2 data are adopted according to recommendations from~\cite{dubey2024llama}, or an equivalent value in bytes. The AdamW optimizer~\citep{adamw} is employed, with $\beta_1 = 0.9$, $\beta_2 = 0.95$, and $\epsilon = 10^{-8}$. The learning rate is scheduled with 2000 steps of linear warm-up, followed by cosine decay to zero. Weight decay is maintained at 0.1, and gradients are clipped globally with a maximum norm of 1.0.

\section{Scaling Trends}
\label{section:scaling}

We offer a comprehensive overview of the scaling behaviors observed in byte-level models, aiming to guide the continued scaling of {\textsc BLT} architectures. Our investigation seeks to overcome constraints noted in prior studies on byte-level models through several strategies: (a) We conduct comparative analyses under compute-optimal training conditions, (b) We develop parallel 8B parameter models trained on substantial datasets (reaching as much as 1T tokens or 4T bytes) and assess performance across downstream tasks, and (c) We analyze scaling properties within fixed inference-cost frameworks. Subsequent sections will explore particular benefits arising from the modeling of byte sequences.

\subsection{Parameter Matched Compute Optimal Scaling Trends}
\label{section:parameter-matched-scaling}
Employing the Llama 2 dataset, we carry out training of multiple \textit{compute-optimal} bpe and {\textsc BLT} models at four distinct parameter scales, spanning from 1B up to 8B parameters. For each configuration, we chart the number of training flops against the language modeling accuracy on a representative sample drawn from the overall training mixture. The bpe models utilize the empirically established optimal proportion of parameters to dataset size, following the methodology described in Llama~3~\citep{dubey2024llama}. This \textit{compute-optimal} configuration is grounded in theoretical results indicating its potential to maximize performance within a restricted training budget~\citep{hoffmann2022training}, thus constituting a rigorous baseline for comparative evaluation. Alongside each bpe model, we train an equivalent {\textsc BLT} model—using an identically sized and architected Latent Transformer—on the same dataset for controlled comparison.

As shown in~\autoref{fig:scaling-compute-opt} (right), {\textsc BLT} models consistently equal or exceed the performance of bpe-based models, with this pattern persisting across increased model scales and flops. To our knowledge, {\textsc BLT} represents the first byte-level Transformer architecture to demonstrate scaling behavior comparable to that of BPE-based models under compute-optimal conditions. This observation supports our hypothesis that the parameter-to-compute ratio that is optimal for bpe models remains effective for {\textsc BLT}, or at least does not significantly deviate.

Both enhancements in model architecture and the adoption of dynamic patching are essential to align with bpe scaling behavior. In ~\autoref{fig:scaling-compute-opt} (left), we present a comparison between models employing space-patching techniques and Llama 3. For SpaceByte \citep{slagle2024spacebyte}, we estimate its performance by implementing {\textsc BLT} space-patching while omitting n-gram embeddings and cross-attention components. While SpaceByte offers gains over Megabyte, it still significantly trails Llama 3. The panel in \autoref{fig:scaling-compute-opt} (right) displays the cumulative benefits achieved through both architectural advancements and dynamic patching approaches. {\textsc BLT} models, when scaled, achieve comparable performance to leading tokenizer-based models like Llama 3.

Furthermore, our results indicate that, for models reliant on tokenizers, the selection of tokenizer impacts overall performance: models utilizing the Llama-3 tokenizer achieve better outcomes than those employing the Llama-2 tokenizer, given identical training datasets.

In summary, the {\textsc BLT} architecture exhibits intermediate behavior between Llama 2 and Llama 3 when operating with much larger patch sizes. While the bpe tokenizers used by Llama 2 and Llama 3 have average token sizes of 3.7 and 4.4 bytes respectively, {\textsc BLT} demonstrates comparable scaling behavior with average patch sizes of 6 and even 8 bytes. Since inference flop count decreases as the average patch size increases, selecting an 8-byte patch size could yield almost a 50\% reduction in inference flops. Furthermore, models configured with larger patch sizes appear to benefit as both model and data scale increases. At the 1B parameter scale, {\textsc BLT} with an 8-byte patch size initially underperforms compared to bpe Llama 2, but it surpasses bpe at the 7B scale. This trend indicates the potential for even larger patch sizes to excel at greater scales, suggesting their viability with continued increases in model size and training computation.

\subsection{Beyond Compute Optimal Task Evaluations}
\label{section:evals}
To further investigate scaling dynamics, we train an 8B {\textsc BLT} model with a data-to-compute ratio that exceeds the established compute-optimal threshold, utilizing the {\textsc BLT}-1T dataset, which offers larger scale and higher data quality. Model performance is then evaluated across a comprehensive array of standard benchmarks for both classification and generation. For these assessments, we focus on tasks assessing common sense reasoning, world knowledge, and code synthesis capabilities:
\paragraph{Classification tasks} assessed include ARC-Easy (0-shot)~\citep{Eval_ARC}, Arc-Challenge (0-shot)~\citep{Eval_ARC}, HellaSwag (0-shot)~\citep{Eval_hellaswag}, PIQA (0-shot)~\citep{Eval_piqa}, and MMLU (5-shot)~\citep{hendrycks2020measuring}. We utilize a prompt-scoring approach, computing the likelihood for each candidate answer, and aggregate results using mean accuracy.
\paragraph{Coding related generation tasks:}  To examine code generation proficiency, we measure pass@1 outcomes on MBPP (3-shot)~\citep{Eval_MBPP} and HumanEval (0-shot)~\citep{Eval_HumanEval}, reflecting the model's capacity to correctly generate Python code on the first attempt.

\autoref{tab:evals} presents a comparison among three models trained on the {\textsc BLT}-1T dataset: a model utilizing a bpe Llama 3 tokenizer,\footnote{We selected the Llama 3 tokenizer, featuring a 128k vocabulary, as it demonstrated superior performance compared to the 32k vocabulary of Llama 2.} alongside two variants of the {\textsc BLT} model. One of these variants adopts a space-patching approach ({\textsc BLT}-Space), while the other implements a patching method based on entropy ({\textsc BLT}-Entropy). The entropy-based model also integrates an approximate monotonicity constraint and resets its context upon encountering new lines, as elaborated in~\autoref{section:entropy-context}. All three systems are trained to equivalent computational (flop) budgets. Notably, for {\textsc BLT}-Entropy, we further adjust the inference-time entropy threshold, lowering it from 0.6 to 0.1. This modification is observed to enhance task performance, though it comes with an increase in the number of inference steps.

The {\textsc BLT}-Entropy model achieves superior results compared to the Llama 3 model on 4 of the 7 evaluated tasks, despite both being trained with an equivalent byte count. This advantage can be attributed to two main factors: (1) enhanced utilization of training compute enabled by dynamic patching, and (2) the explicit incorporation of byte-level information rather than relying solely on token representations.

In contrast, {\textsc BLT}-Space performs worse than the Llama 3 tokenizer on all tasks except one. Nevertheless, it notably lowers inference flops due to its increased average patch size of 6 bytes. By comparison, both the bpe and entropy-patching-based approaches achieve similar average patch sizes of about 4.5 bytes within the same training dataset. Given an identical training budget, the model with the larger patch size is able to process 30\% more data than the other two approaches, which could result in {\textsc BLT} deviating further from the point of computational optimality.

\subsection{Patches Scale Better Than Tokens} {\textsc BLT} models enable the concurrent scaling of both model size and patch size, all while preserving a constant budget for training and inference computations and using the same amount of training data. Unlike traditional token-based models with a fixed vocabulary, patch-based models uniquely allow for patch size to be increased as desired—an advantage that circumvents the efficiency limitations highlighted in Section~\ref{section:bpe-patch}. Employing longer patch sizes reduces the computation required per forward pass, which in turn frees resources that can be allocated to expand the global latent transformer, since it is executed less frequently.

We perform a fixed-inference scaling analysis to examine whether increasing model size while reducing the number of steps on larger patches yields superior performance compared to smaller models executing more steps. Using initial model sizes of 400m and 3.6B parameters with the Llama 2 tokenizer, we identify models with equivalent FLOP counts using the Llama 3 tokenizer and {\textsc BLT}-Entropy architectures employing average patch sizes of 6 and 8 bytes within the training data mix (refer to~\autoref{tab:fixed-inf-params} for specific model configurations). For models with a patch size of 8, we opt for 3 encoder layers rather than just 1. Each configuration is trained across several training FLOP allocations.

As illustrated in \autoref{fig:fixed_inference_scaling}, {\textsc BLT} models demonstrate superior scaling performance relative to tokenization-based models across both inference flop regimes. Notably, with limited training resources, BPE models initially outperform, but are soon overtaken by {\textsc BLT} models as training scales slightly beyond the compute-optimal threshold. In many scenarios, allocating additional resources to one-time pretraining is advantageous, as it yields models with enhanced performance within a predetermined inference cost. A clear illustration is seen in the case of 8B models such as Llama 3.1, which has been exposed to approximately two orders of magnitude more data than the compute-optimal amount for models of its size.

The point at which {\textsc BLT} surpasses token-based models in performance has moved incrementally closer to the compute-optimal configuration as model size increases within the larger flop class, shifting from 3x down to 2.5x the optimal compute budget. Furthermore, the model utilizing patch size 8 displays a steeper improvement trajectory in this high-flop regime, surpassing alternative models at an earlier stage. As outlined in Section~\ref{section:parameter-matched-scaling}, models with larger patch sizes tend to approach the performance of BPE models at increased parameter counts. This phenomenon can be partly explained by the diminishing proportion of computational resources allocated to the byte-level Encoder and Decoder components, whose scaling is more modest compared to that of the Latent Transformer. When expanding the parameter count twentyfold, from 400M to 8B, the local parameters of {\textsc BLT} only increase by a factor of approximately two. This is notable because larger patch sizes exclusively impact the computational demands of the patch Latent Transformer, leaving the byte-level modules largely unaffected. Consequently, when progressing to higher model scales, the {\textsc BLT}-Entropy model with ps=8 grew from 1.6x to 1.7x the size of the Llama 2 model, reflecting this dynamic.

To summarize, our investigation into patch-length scaling reveals that the {\textsc BLT} patch-based architecture benefits from improved scaling performance when both the patch size and the overall model size are increased together. Notably, this improvement persists and can even be accentuated as models become larger.

\section{Byte Modeling Improves Robustness} We further assess the robustness of {\textsc BLT} relative to token-based architectures that do not have inherent access to byte-level data, and introduce a method for augmenting pretrained token-based models to handle byte-level information.
\label{sec:robustness}

\subsection{Character-Level Tasks}

An initial impetus behind the development of byte-level models was their demonstrated resilience to noise at the byte level within input data, coupled with their inherent sensitivity to the underlying structure of tokens—a challenge for existing models reliant on tokenizers. To assess these capabilities, we conduct supplementary experiments using benchmark datasets that specifically test resistance to input noise and examine the models' sensitivity to individual characters (including not only those in English, but also in multiple languages, as well as numeric digits and phonemes). The outcomes of these evaluations are shown in Table~\ref{tab:char_tasks}.

\paragraph{Noisy Data} We generate altered versions of the standard benchmark classification datasets outlined in Section~\ref{section:evals} in order to assess and compare the robustness of {\textsc BLT} against tokenizer-based models. To induce textual variation, we utilize five unique character-level noise techniques: (a) \textit{AntSpeak}, which transforms all characters to uppercase and separates them with spaces; (b) \textit{Drop}, which deletes 10\% of the characters in the text at random; (c) \textit{RandomCase}, which assigns uppercase or lowercase to 50\% of the characters each, randomly distributed throughout the text; (d) \textit{Repeat}, where 20\% of the characters are repeated up to four consecutive times; and (e) \textit{UpperCase}, converting the entire text to uppercase. In the evaluation process, these noising strategies are applied independently to the prompt, the completion, or to both components, resulting in separate evaluation tasks. Average scores are reported accordingly. Table~\ref{tab:char_tasks} summarizes the performance on noised HellaSwag~\citep{Eval_hellaswag}, demonstrating that {\textsc BLT} consistently achieves higher robustness compared to tokenizer-based counterparts, surpassing them by an average of 8 points, even when compared to models trained on the same data, and outperforming the more extensively trained Llama 3.1 model as well.

\paragraph{Phonology - Grapheme-to-Phoneme (G2P)} We evaluate how effectively {\textsc BLT} converts sequences of graphemes (the characters constituting a word) into phonemic representations, which correspond to the word’s pronunciation. The results for the G2P task, conducted within a 5-shot framework and utilizing Phonology Bench~\citep{suvarna2024phonologybench}, are reported in Table \ref{tab:char_tasks}. Our findings indicate that {\textsc BLT} achieves superior performance compared to the baseline Llama 3 1T tokenizer-based model on this task.

\paragraph{CUTE}
To evaluate character-level comprehension, we test {\textsc BLT} using the CUTE benchmark~\citep{edman2024cute}, which includes multiple tasks grouped into three primary types: compositional understanding, orthographic similarity, and sequence manipulation. This benchmark is notably difficult for most models that rely on tokenization; these models tend to recognize token spellings yet have difficulty applying this knowledge to perform text manipulation. As indicated in Table~\ref{tab:char_tasks}, {\textsc BLT}-Entropy surpasses both BPE Llama 3 variants on this benchmark by a margin exceeding 25 points. Notably, our model excels in character manipulation tasks, attaining 99.9\% accuracy on both spelling tasks. These substantial gains occur even though {\textsc BLT} has been trained with 16 times less data than Llama 3.1, underscoring the challenge BPE-based models face when acquiring character-level representations. Figure~\ref{fig:cute_examples} presents examples in which the Llama 3 tokenizer-based model underperforms relative to our {\textsc BLT} model. The only tasks in which BPE is superior are word deletion and insertion; while such word-based operations may pose additional challenges for byte-level architectures, the performance difference is modest, suggesting that building from character-level to word-level operations may ultimately be more tractable than the reverse. Evaluation was carried out uniformly across tasks using the original Huggingface prompts. It is possible that prompt engineering could further enhance BPE model performance.

\paragraph{Low Resource Machine Translation} We assess the performance of {\textsc BLT} on translation tasks involving both directions for six major language families and twenty-one low-resource languages, which feature a range of scripts, utilizing the FLORES-101 benchmark~\citep{goyal2022flores}. SentencePiece BLEU scores are reported in Table~\ref{tab:flores}. The experimental findings reveal that {\textsc BLT} surpasses the model utilizing the Llama 3 tokenizer, yielding an average improvement of 2 BLEU points for translations into English and a 0.5-point boost for translations from English. While performance on widely-used language pairs is generally on par with or marginally exceeds that of Llama 3, {\textsc BLT} delivers notably better results across many lower-resource language pairs. These results highlight the strength of byte modeling for capturing and generalizing over extended and varied byte sequences in low-resource settings.

\subsection{Training {\textsc BLT} from Llama 3}
\label{sec:distillation}
We investigate a method where {\textsc BLT} models can take advantage of existing pre-trained tokenizer-based architectures to achieve improved and more efficient training. This is accomplished by initializing the global transformer parameters of {\textsc BLT} with those obtained from a pre-trained Llama 3.1 model. After initialization, the global transformer parameters are updated using a learning rate set to one-tenth of that used for the local encoder and decoder components, following the optimal training schedule prescribed for Llama 3. A performance comparison with a baseline {\textsc BLT} model is reported in Table~\ref{tab:distill}. The results clearly indicate that the {\textsc BLT} model initialized from Llama 3.1 substantially outperforms both the Llama 3 baseline and the baseline {\textsc BLT}, even though all models were trained with an equivalent compute budget (in flops). Furthermore, relative to our {\textsc BLT}-Entropy model—described in Table~\ref{tab:evals}—which was trained on a much larger corpus (1T tokens), the {\textsc BLT} initialized from Llama 3.1 still achieves higher accuracy on the MMLU benchmark. This demonstrates that such an approach can be a highly effective strategy for substantially decreasing training flops while maintaining or even improving downstream performance.

This approach may be interpreted as adapting models that rely on tokenizers into those that function without them, essentially transforming a pre-trained LLaMA 3.1 model into a {\textsc BLT} system. For a thorough evaluation, we present the baseline LLaMA 3.1, trained on 15T tokens, in Table \ref{tab:distill}, and contrast its outcomes with those of the {\textsc BLT} model derived from LLaMA 3. While our method produces only marginally lower results on MMLU and HumanEval, it exhibits noticeably greater drops in performance across several other benchmarks. These results highlight that further research is necessary to maximize the utility of the pre-trained model, particularly by fine-tuning aspects such as data mixture strategies and various hyperparameters.

\section{Ablations and Discussion}
\label{section:ablations}
Here, we present ablation studies that clarify the motivations behind the design decisions for {\textsc BLT}, including choices regarding the patching approach and selection of hyper-parameters for the {\textsc BLT} model with 8B parameters, as trained on the {\textsc BLT}-1T dataset.

\paragraph{Entropy Model Hyper-parameters} To investigate how the size of the entropy model and the length of the context window influence scaling dynamics, we develop a series of byte-level entropy transformer architectures spanning from 1 million to 100 million parameters, systematically varying the context window lengths between 64 and 512 tokens. Scaling law plots depicting bits per byte (bpb) against the number of training flops are generated, utilizing our $400m$ and $1b$ {\textsc BLT} models trained with the Llama-2 dataset; these results are illustrated in \autoref{fig:entropy_ablation}. Our analysis reveals that increases in both entropy model size and context window length generally enhance scaling behavior, though gains become marginal past approximately 50 million parameters.

\paragraph{Types of Patching}

We conduct an ablation study on the four distinct patching strategies outlined in Section~\ref{section:patching}: 1) Strided Patching employing strides of 4 and 6, 2) Whitespace-based Patching, 3) Patching determined by a BPE Tokenizer following the Llama 3 tokenizer, and 4) Entropy-driven patching leveraging a lightweight byte-level language model.

Although dynamic patching effectively shortens sequence lengths, we regulate sequence length to ensure that every patching approach preserves an equivalent context length. Consequently, each model, on average, is exposed to an identical quantity of bytes per sequence during both training and inference. This eliminates potential confounds arising from the ability to process longer contexts. Results from these ablation studies are depicted in Figure~\ref{fig:scaling-compute-opt}. Notably, all alternative patching schemes demonstrate superior performance relative to static patching, with space patching showing nearly comparable effectiveness to dynamic entropy based patching.

\autoref{tab:evals_ablation} displays comparative benchmark results for {\textsc BLT} models utilizing different patching strategies—namely tokenizer-based, space patching, and entropy-based patching—each trained on the Llama 2 dataset for an optimized number of training steps~\citep{dubey2024llama}. While space patching offers a more straightforward approach by avoiding the computational overhead of running an entropy model during training, our results demonstrate that the improvements afforded by entropy-based patching, as discussed in Section~\ref{section:scaling}, also extend to a range of downstream benchmark evaluations.\footnote{The outcomes for space patching are based on earlier experiments conducted without cross-attention, though parallel patterns persist when cross-attention is applied.}

\paragraph{Cross-Attention}
As shown in~\autoref{tab:crossattn_ablations_1b}, we investigate the impact of introducing cross-attention at different locations within both the encoder and decoder of {\textsc BLT}. For encoder cross-attention, our experiments consider three query initialization strategies: (1) a single shared learned embedding applied to all global states, (2) a hash embedding generated from the bytes present in the patch, and (3) pooling operations over the encoder’s hidden states for the patch bytes at the corresponding layer.

Our results indicate that implementing cross-attention within the \textit{decoder} yields the strongest performance gains. When applied to the encoder, cross-attention offers only marginal benefits, and these are confined to scenarios where query initialization is accomplished through pooling. Furthermore, our analysis reveals that cross-attention provides noticeable advantages for the Common-Crawl dataset, with the positive impact being more pronounced as patch sizes increase.

\paragraph{n-gram Hash Embeddings} 

We examine the effect of varying n-gram hash embedding vocabulary sizes—specifically 0, 100K, 200K, and 400K—and summarize our findings in Table~\ref{tab:ngram_ablations_1b}. Our analysis reveals that hash embeddings confer benefits across every domain considered, with the most pronounced improvements on Wikipedia and Github datasets (a 0.04 bpb improvement versus 0.01 bpb at 15k steps with the 8B model). Scaling the model to 8B and reducing the number of hash bins from 500K to 300K only led to a marginal ~0.001 bpb drop at 15k steps. These results demonstrate the essential role of hash embeddings in elevating {\textsc BLT} models to reach the performance levels of their tokenizer-based counterparts; nevertheless, beyond 300K hash bins, further gains become minimal. Moreover, our observations suggest that these performance increases are largely independent from those achieved through cross-attention mechanisms, as both approaches offer advantages on distinct datasets.

\paragraph{Local Model Hyperparameters} Table \ref{tab:local_layers} presents an ablation study examining different configurations for the number of layers in both the local encoder and decoder. Results indicate that, in combination with hash n-gram embeddings, {\textsc BLT} achieves favorable performance when the encoder is kept very shallow—with only a single layer—while the decoder utilizes a deeper architecture.

\section{Related Work}
\label{section:related_work}

\textbf{Character-Level RNNs:} The task of Character Language Modeling has long been established as fundamental within the neural model community~\citep{sutskever2011generating,mikolov2012subword,graves2013generating}, largely due to its innate capacity to seamlessly handle out-of-vocabulary items, circumventing the need for traditional back-off strategies. \cite{kim2016character} introduced an architecture that exclusively processes characters as input, employing convolutional and highway layers as a prelude to LSTM-based RNNs. Their approach achieves parity with then-state-of-the-art RNN language models on English and attains superior outcomes for morphologically complex languages—an often-noted benefit of character-level LLMs. Expanding this line of inquiry, \cite{kenter2018byte} conducted machine comprehension experiments with byte-level LSTM architectures, once again showing that such models surpass their word-based counterparts for morphologically rich languages like Turkish and Russian. In related research, \cite{zhang2015character} implemented character-level convolutional neural networks for text classification tasks, demonstrating improvements over word-level baselines in specific cases. Building on these methodologies, \citet{chung2019hierarchical} adopted hierarchical LSTM systems with embedded boundary detectors at each tier to automatically extract text hierarchies, further elevating character-level modeling performance. ByteNet, proposed by \cite{kalchbrenner2016neural}, utilizes convolutional neural network layers at the character level, providing an alternative to attention mechanisms in the context of machine translation.

\textbf{Character-Level Transformers:} Transformer architectures, which leverage attention mechanisms~\citep{vaswani2017attention} and subword tokenization strategies~\citep{sennrich-etal-2016-neural}, have markedly advanced neural models' capabilities on various language modeling benchmarks. Nevertheless, employing word or sub-word units introduces an inherent inductive bias regarding the desired level of abstraction for these models. Seeking to bridge transformer methodologies with earlier promising findings from character-level language modeling, \citet{al2019character} deployed highly deep transformer networks. By integrating auxiliary loss functions, they trained transformer-based character models that surpassed the performance of prior character-level LSTM language models. Despite these improvements, a substantial performance difference remained relative to word-level large language models. Additionally, GPT-2~\citep{radfordgpt2} reported that byte-level language models failed to reach the effectiveness of their word-level counterparts on expansive corpora such as the 1 billion word benchmark.

Although \cite{Choe2019BridgingTG} showed that transformer-based byte-level LLMs can surpass their subword-based counterparts with similar parameter counts, these models require substantially higher computational resources and extended training times. Along similar lines, \cite{el2020characterbert} introduced CharFormer, a BERT variant that constructs word embeddings by processing character-level representations via convolutions. While CharFormer demonstrated superior performance in the medical domain, it came at the cost of increased computational expenditure. In another effort, \cite{clark2022canine} proposed CANINE, an encoder-only model with 150M parameters, operating directly on character sequences. CANINE's architecture features a deep transformer stack analogous to our global model, augmented with a local transformer and strided convolution layers for input character downsampling. This design led CANINE to exceed the performance of the comparably sized token-based encoder-only model, mBERT, across multiple multilingual downstream tasks. Meanwhile, ByT5~\citep{xue2022byt5} investigated byte-level encoder-decoder models that forego patching operations entirely. Though these models demonstrated stronger resilience to noisy inputs and matched the performance of tokenized models when trained with just a quarter of the data, the absence of patching resulted in costly attention calculations across every byte—making the approach computationally intensive. Modeling inputs at the byte rather than subword unit level leads to longer input sequences, which poses challenges to the scalability of such architectures. More recently, \cite{wang2024mambabyte} leveraged the Mamba Architecture~\citep{gu2023mamba}, characterized by its ability to uphold a fixed-size memory state across very large context windows, to construct a byte-level Mamba-based model without using patching. At the 350M parameter level and under compute-matched conditions, their approach demonstrated improved bits-per-byte performance over byte-level transformer models on several benchmarks.

\textbf{Patching-based approaches:} Employing patching techniques has been shown to reduce the otherwise high computational demands (in terms of flops) of byte-level LLMs, sometimes without substantial loss in model effectiveness. Several studies have documented early successes at modest scales in both model capacity and data volume. For example, \cite{nawrot-etal-2022-hierarchical} utilized static patching for downsampling and upsampling in their design of the hourglass transformer, which exceeded the performance of alternative byte-level baselines at the 150M parameter level. Building upon this, \cite{nawrot-etal-2023-efficient} introduced improvements via dynamic patching strategies, proposing a boundary-predictor trainable in an end-to-end setting, a boundary-predictor guided by specific tokenizers, and an entropy-driven patching method inspired by {\textsc BLT}. Their findings indicate that these advances can surpass contemporaneous vanilla transformers on language modeling benchmarks, particularly for models with 40M parameters trained on approximately 400M tokens. Meanwhile, \cite{lester2024training} explored models trained on data sequences compressed using arithmetic coding, obtaining higher compression ratios than those achievable with BPE. By implementing an equal-info windows strategy, they were able to beat byte-level models on language modeling benchmarks, though their performance fell short of leading subword-based baselines.

This research takes significant cues from the design of MegaByte~\citep{yu2023megabyte}, a decoder-only causal LLM that leverages a predetermined, static scheme to convert bytes into fixed-sized patches by concatenating byte representations. MegaByte further employs a local model within the decoder to reconstruct bytes from these patches. The original study demonstrates that MegaByte can perform on par with traditional tokenizer-based models when scaled to 1B parameters, evaluated on a dataset encompassing approximately 400B bytes. In our study, we include comprehensive ablation of MegaByte and find that static patching consistently falls short relative to the latest tokenizer-based models trained optimally for computational efficiency under FLOP constraints; we show that {\textsc BLT} effectively addresses and narrows this performance gap. Notably, \cite{slagle2024spacebyte} independently report similar findings about MegaByte’s shortcomings. They propose extending the static patching approach to operate at whitespaces and other whitespace-like byte boundaries, complemented with the addition of a local encoder; their findings indicate performance gains over tokenized transformer baselines on certain datasets such as Github and arXiv at the 1B parameter scale in compute-controlled comparisons. We include new experiments with this variant, concluding that further architectural enhancements are essential for byte-level models to fully scale and truly rival the strongest token-based models, such as Llama 3.

\section{Limitations and Future Work}

In this study, when making architectural decisions, we train our models for the number of steps identified as optimal for Llama~3~\citep{dubey2024llama}. Nonetheless, these scaling laws were derived with BPE-level transformers in mind, and may result in less than ideal (data, parameter size) proportions when applied to {\textsc BLT}. Deriving accurate scaling laws tailored to {\textsc BLT} is a direction we reserve for future research, as this may yield even more advantageous scaling characteristics for our model. Furthermore, most of our experiments are limited to model sizes up to 1B parameters, so there is a possibility that the most effective architectural choices could shift as we increase scale to 8B parameters and higher, potentially allowing for superior performance at larger model sizes.

Current transformer codebases and libraries primarily focus on optimizing efficiency for transformer models that operate using tokenizers. Although our study includes theoretical flop-matched experiments and leverages some specialized implementations—like FlexAttention—for handling components that differ from standard transformer layers, our models may still lag behind tokenizer-based counterparts in wall-clock performance. Additional optimization efforts could thus further improve their runtime efficiency.

Whereas {\textsc BLT} employs an entropy model for patching that is trained independently, exploring approaches to jointly train the patching model end-to-end represents a promising avenue for upcoming studies. Section \ref{sec:distillation} introduces preliminary experiments that provide early evidence supporting the feasibility of ``byte-ifying'' large tokenizer-based architectures, like Llama 3, which have been trained on datasets exceeding 10T tokens. This is achieved by initializing the global transformer with the pretrained weights and then keeping those weights fixed. Further investigation along these lines could reveal strategies that not only preserve the advantages introduced by bytefying but also enhance performance relative to the original tokenizer-based models, all without necessitating complete retraining.

\section{Conclusion}
\label{section:conclusion}

In this work, we introduce the Byte Latent Transformer (\textbf{BLT}), an innovative architecture that challenges the standard reliance on fixed-vocabulary tokenization in large language models. {\textsc BLT} proposes a flexible, trainable approach to aggregating bytes into patches, thereby enabling the model to distribute computational effort according to the intrinsic complexity of the input data. This adaptive strategy yields substantial gains in both computational efficiency and model robustness. Through comprehensive scaling experiments, we show that {\textsc BLT} achieves parity with tokenization-based models such as Llama 3 across scales reaching 8B parameters and 4T bytes of data, and is capable of reducing inference flops by as much as 50\% with only minor reductions in standard evaluation metrics. Additionally, {\textsc BLT} introduces a novel scaling avenue, permitting concurrent increases in both model capacity and patch granularity within the same inference computational limits—a benefit particularly relevant for practical deployment scenarios. By processing inputs directly at the byte level, {\textsc BLT} also demonstrates enhanced handling of rare and atypical data, yielding notable improvements in resilience to noisy samples and offering richer representations of sub-word information. Collectively, these findings establish {\textsc BLT} as a compelling alternative to conventional tokenization-based architectures, presenting a robust, scalable foundation for more flexible and efficient language modeling.

\end{document}